% chapters/08/sub-cpmk.tex
\begin{subcpmkbox}
Bab ini adalah bab pertama pada \CPMK{2}, yaitu penerapan algoritma kecerdasan
buatan dasar untuk NPC. Bab ini menuntaskan \SUBCPMK{2.1} secara utuh.

\vspace{0.5em}
\textbf{\SUBCPMK{2.1} Merancang perilaku NPC sederhana (misalnya
\istilah{chase}) berdasarkan posisi pemain} --- ranah C3, \CPMK{2},
\CPL{03}.

\vspace{0.5em}
{\setlength{\tabcolsep}{4pt}\footnotesize%
\begin{tabularx}{\textwidth}{@{}C{1.9cm} L{4.0cm} L{2.5cm} Y@{}}
\toprule
\headrow \textbf{Indikator} & \textbf{Rumusan} & \textbf{Dibahas pada} & \textbf{Bukti} \\
\midrule
\ind{a} & Menghitung vektor arah NPC menuju pemain dan menormalkannya &
Bagian~\ref{sec:persepsi} dan Bagian~\ref{sec:menormalkan} &
Kode perhitungan arah pada kelas NPC, beserta penjelasan setiap langkahnya. \\
\addlinespace
\ind{b} & Membatasi besar kecepatan agar gerak NPC konsisten &
Bagian~\ref{sec:menormalkan} &
Tabel pengukuran kecepatan pada berbagai jarak, yang menunjukkan kecepatan
tetap, beserta perbandingannya dengan cara yang salah. \\
\addlinespace
\ind{c} & Menjelaskan perbedaan AI reaktif dan AI berbasis perencanaan &
Bagian~\ref{sec:reaktif-perencanaan} &
Jawaban soal uraian, beserta catatan pada lembar kerja praktikum mengenai
kelemahan perilaku reaktif yang Anda amati sendiri. \\
\bottomrule
\end{tabularx}}
\end{subcpmkbox}
